\documentclass[aspectratio=169,11pt,t]{beamer}

% ---------- Style: matches the Week 2, 3 and 4 practice talk decks ----------
\usepackage[T1]{fontenc}
\usepackage{mathpazo}
\usepackage{helvet}
\usepackage{graphicx}
\usepackage{tikz}

\definecolor{navy}{HTML}{1F3864}
\definecolor{gold}{HTML}{F2C14E}
\definecolor{midgray}{HTML}{595959}
\definecolor{hairline}{HTML}{C8C8C8}
\definecolor{palegray}{HTML}{BFBFBF}
\definecolor{midblue}{HTML}{6E8FBF}
\definecolor{teal}{HTML}{3E7C7B}
\definecolor{grayline}{HTML}{8C8C8C}

\setbeamertemplate{navigation symbols}{}
\setbeamertemplate{itemize items}{\textcolor{navy}{\footnotesize\textbullet}}
\setbeamercolor{normal text}{fg=black}
\setbeamercolor{structure}{fg=navy}
\setbeamerfont{normal text}{family=\rmfamily}
\usefonttheme{serif}

% Week 4 used a topic label here. Week 5's subject is the assertion headline, so the
% frame title states the finding as a full sentence instead of naming the subject.
\newcommand{\assertion}[1]{%
  {\sffamily\bfseries\fontsize{15}{19}\selectfont\color{navy}#1}\par
  \vspace{4pt}\textcolor{hairline}{\rule{\textwidth}{0.5pt}}\par\vspace{10pt}}

\begin{document}

% ---------- Slide 1: Title ----------
% Says who is speaking and what the talk is about, and nothing else. The subtitle is
% already the assertion the whole talk defends.
\begin{frame}
  \vspace{55pt}
  {\sffamily\fontsize{8}{11}\selectfont\color{midgray}MSDS 696 DATA SCIENCE PRACTICUM II \quad\textbar\quad WEEK 5}\par\vspace{12pt}
  {\sffamily\bfseries\fontsize{25}{30}\selectfont\color{navy}Predicting U.S. Bank Distress}\par\vspace{9pt}
  {\fontsize{12.5}{17}\selectfont\itshape\color{midgray}A proposal for ranking which banks examiners review first}\par\vspace{14pt}
  {\sffamily\fontsize{9.5}{13}\selectfont Oussama Ennaciri \quad\textbar\quad August 2026}\par\vspace{6pt}
  \textcolor{navy}{\rule{\textwidth}{1.2pt}}
\end{frame}

% ---------- Slide 2: The visual ----------
% The one visual the assignment asks for. The headline claims the model ranks better,
% so the evidence has to be every threshold rather than one: a single-budget bar chart
% would only support the claim at one point. Both axes start at zero.
% Vertical space is set explicitly rather than with \vfill: \vfill pushed the headline
% flush to the top edge and the axis label flush to the bottom, so the slide read as
% having no margins at all.
\begin{frame}
  \vspace*{10pt}
  \assertion{The model ranks risky banks better than the currently used benchmark}
  \vspace{4pt}
  \centerline{\includegraphics[width=0.93\textwidth]{../figures/final_2_gaincurve_slide.png}}
\end{frame}

\end{document}
